\begin{resumo}

\textcolor{black}{
Este trabalho apresenta o desenvolvimento de um micromouse autônomo projetado para navegar de forma eficiente em labirintos de diferentes dimensões, buscando alcançar uma área central específica sem qualquer intervenção humana após sua ativação inicial. O sistema deve executar, de maneira independente, tarefas fundamentais como detecção de paredes, monitoramento contínuo da própria posição, mapeamento progressivo do ambiente e identificação do ponto que caracteriza o objetivo do desafio. Para garantir a padronização, foram definidas especificações técnicas que incluem limites físicos do robô — que não deve ultrapassar 25 centímetros em comprimento ou largura — e sua capacidade de adaptação a três configurações distintas de labirinto, variando entre 4x4, 8x8 e 16x16 células, todas estruturadas com dimensões e padrões visuais específicos. Diante dessas condições, o micromouse precisa interpretar corretamente o ambiente, ajustando suas decisões de navegação de acordo com a disposição das paredes e com as possibilidades de rota até o destino final. O projeto também incorpora um sistema de telemetria capaz de exibir em tempo real dados como trajetória percorrida, consumo de bateria, velocidade média e tempo total de execução. Após cada tentativa, essas informações são registradas em banco de dados, permitindo análises posteriores e a comparação de desempenhos em diferentes labirintos, promovendo uma compreensão mais ampla do comportamento do sistema autônomo desenvolvido. Dessa forma, o projeto proporciona uma experiência prática e multidisciplinar que envolve diversas engenharias, estimulando a aplicação integrada desses conceitos e contribuindo para a formação técnica no desenvolvimento de soluções de navegação inteligente em ambientes estruturados.
}

\vspace{\onelineskip}
\noindent
\textbf{Palavras-chave}: \textcolor{red}{\imprimirpalavrachaveum; \imprimirpalavrachavedois.}

\end{resumo}

% \begin{resumo}
%  O resumo deve ressaltar o objetivo, o método, os resultados e as conclusões 
%  do documento. A ordem e a extensão
%  destes itens dependem do tipo de resumo (informativo ou indicativo) e do
%  tratamento que cada item recebe no documento original. O resumo deve ser
%  precedido da referência do documento, com exceção do resumo inserido no
%  próprio documento. (\ldots) As palavras-chave devem figurar logo abaixo do
%  resumo, antecedidas da expressão Palavras-chave:, separadas entre si por
%  ponto e finalizadas também por ponto. O texto pode conter no mínimo 150 e 
%  no máximo 500 palavras, é aconselhável que sejam utilizadas 200 palavras. 
%  E não se separa o texto do resumo em parágrafos.

%  \vspace{\onelineskip}
    
%  \noindent
%  \textbf{Palavras-chave}: latex. abntex. editoração de texto.
% \end{resumo}

% O resumo é um item obrigatório. Essa parte do relatório será uma visão rápida e clara do projeto desenvolvido. O leitor terá informações como: descrição breve do projeto, principais requisitos, tecnologias necessárias e outras informações relevantes para apresentação. O resumo terá no máximo meia (1/2) página.}

% \textcolor{red}{Ao longo deste texto, as descrições em \textbf{cor vermelha} são meras instruções, não devendo aparecer na versão final do texto. \textbf{Utilize o comentário \% ao invés de apagar estas descrições, para não perder as orientações apresentadas.}}